% Generated draft: requires mathematical and source review.
\documentclass[11pt]{article}
\usepackage[margin=25mm]{geometry}
\usepackage{fontspec}
\setmainfont{Latin Modern Roman}
\usepackage{amsmath,amssymb,mathtools,unicode-math}
\setmathfont{Latin Modern Math}
\usepackage{enumitem,needspace,xurl,hyperref}
\setlength{\parindent}{0pt}
\setlength{\parskip}{5pt}
\setlength{\emergencystretch}{3em}
\newcommand{\N}{\mathbb{N}}
\newcommand{\Z}{\mathbb{Z}}
\newcommand{\Q}{\mathbb{Q}}
\newcommand{\R}{\mathbb{R}}
\newcommand{\C}{\mathbb{C}}
\newcommand{\F}{\mathbb{F}}
\newcommand{\A}{\mathbb{A}}
\newcommand{\PP}{\mathbb{P}}
\newcommand{\E}{\mathbb{E}}
\newcommand{\eps}{\varepsilon}
\newcommand{\dd}{\,\mathrm{d}}
\begin{document}
\section*{UnsolvedMath ID: 30005081}
\textbf{Source problem code:} catalogue entry
\subsection{Definitions and mathematical statement}
Let \(X\) be a projective variety over a field \(k\). The supplied statement does not specify the base field or the permitted class of varieties beyond projectivity.

For terminology only, under the normal split convention, a toric variety over \(k\) is a normal integral variety \(Y\) containing a dense open torus \(T\cong(\mathbb G_{m,k})^r\), whose multiplication action on itself extends to an action on \(Y\). One common model of a toric degeneration of \(X\) is a flat projective morphism
\[
\pi\colon\mathcal X\longrightarrow\mathbb A^1_k
\]
whose fiber over \(0\) is a toric variety and for which there is a nonempty open subset \(U\subseteq\mathbb A^1_k\setminus\{0\}\) such that every geometric fiber over \(U\) is isomorphic to the corresponding base change of \(X\). This is an illustrative convention: the supplied record does not determine the parameter space, normality or splitness requirements, or the precise meaning of the general fiber being \(X\).

\textbf{Research problem.} For a projective variety \(X\), develop general methods to compute toric degenerations of \(X\), and characterize the relationships among distinct toric degenerations of the same \(X\).

The source presents two methodological questions rather than a quantified conjecture. It does not specify an effective presentation of \(X\), the required output of a computation, an equivalence relation on degenerations, or which relationships are to be characterized. Consequently, the requested methods and comparison problem require further specification before they become a single precise algorithmic or classification statement. The original excerpt refers more broadly to an algebraic variety; projectivity is retained from the supplied reconstructed statement.
\subsection{Short English statement}
The problem asks how to construct toric degenerations of a projective variety and how to compare different degenerations of that variety. The base field, admissible families, toric conventions, and intended comparison relation remain unspecified. This is a statement-only draft requiring human mathematical review.
\subsection{Sources}
\begin{itemize}
\item Supplied local UnsolvedMath registry record, permanent ID \(30005081\), title: \emph{Computation and Comparison of Toric Degenerations}; exact source problem code: \texttt{catalogue entry}.
\item Supplied source citation: \emph{Toric Geometry (2023)}, Oberwolfach Reports, DOI \texttt{10.4171/owr/2022/17}, pp.\ 861--927. \url{https://doi.org/10.4171/owr/2022/17}. The original excerpt attributes \emph{Computing toric degenerations of varieties} to Fatemeh Mohammadi, joint work with Oliver Clarke and Francesca Zaffalon. The registry records proposed year \(2022\), with no separately named proposer.
\item Supplied literature source: Clarke--Mohammadi--Zaffalon. \url{https://doi.org/10.1016/j.jalgebra.2023.09.026}.
\item Supplied literature source: Bossinger--Mohammadi--Nájera Chávez. \url{https://doi.org/10.3842/SIGMA.2021.059}.
\item Supplied literature source: Oberwolfach Report 2022/17. \url{https://doi.org/10.4171/owr/2022/17}.
\item Literature assessment dated exactly \texttt{2026-08-22}: supplied dated curation, not independently verified. The registry labels the problem partially solved and reports that Gröbner, tropical, and matching-field methods compute many toric degenerations and relate important families, while no general method or classification of all degenerations and their relationships is known. The background identifies this as literature triage derived from the supplied Oberwolfach curation and requests independent verification of time-sensitive status and cited claims.
\item The supplied statement status is \texttt{corrected\_verified}; its verification field says that the statement was reconstructed and checked against the official Oberwolfach report. These are supplied provenance claims, not independent verification performed for this draft.
\end{itemize}
\end{document}
